\chapter{Procesy gaussowskie: od kowariancji do prawa procesu}
\label{ch:gauss}

\section{Motywacja: czy średnia i kowariancja opisują cały proces?}
W rozdziale pierwszym zobaczyliśmy, że znajomość rozkładu $X_t$ dla
każdego $t$ nie wystarcza do określenia prawa procesu: potrzebujemy
rozkładów \emph{łącznych} dla wszystkich skończonych zbiorów chwil.
W ogólności obliczanie tych rozkładów jest bardzo trudne. Procesy
gaussowskie są niezwykłe właśnie dlatego, że wszystkie ich rozkłady
skończeniewymiarowe można odtworzyć z dwóch funkcji:
$m(t)=\E X_t$ oraz $K(s,t)=\Cov(X_s,X_t)$.

Warto od razu oddzielić \emph{gaussowskość pojedynczych zmiennych}
od \emph{gaussowskości całego procesu}. Można mieć $X_s\sim\Normal(0,1)$
i $X_t\sim\Normal(0,1)$, choć para $(X_s,X_t)$ nie ma rozkładu
normalnego dwuwymiarowego. Dopiero odpowiedni warunek na każdą
skończoną rodzinę rozstrzyga o gaussowskości procesu.

\begin{definicja}[Wektor gaussowski]
Niech $d\in\N$ i $Z=(Z_1,\ldots,Z_d)$ będzie wektorem losowym
o wartościach w $\R^d$. Nazywamy go \emph{wektorem gaussowskim},
jeżeli dla każdego $a\in\R^d$ zmienna $\sum_{j=1}^d a_jZ_j$
ma rozkład normalny (dopuszczamy rozkład zdegenerowany, czyli
skoncentrowany w jednym punkcie). Jeżeli $m=\E Z$ oraz
$\Sigma=(\Cov(Z_i,Z_j))_{i,j=1}^d$, zapisujemy
$Z\sim\Normal_d(m,\Sigma)$.
\end{definicja}

Dopuszczenie degeneracji nie jest tylko wygodą techniczną. Dla
$Z=(U,U)$, gdzie $U\sim\Normal(0,1)$, wektor $Z$ jest gaussowski,
ale jego rozkład jest skoncentrowany na przekątnej $\{(x,x):x\in\R\}$,
więc nie ma gęstości względem dwuwymiarowej miary Lebesgue'a.

\begin{stwierdzenie}[Charakterystyka wektora gaussowskiego]
Dla każdego $d\ge1$, $m\in\R^d$ i symetrycznej nieujemnie
określonej macierzy $\Sigma\in\R^{d\times d}$ istnieje wektor
$Z\sim\Normal_d(m,\Sigma)$, którego funkcja charakterystyczna ma postać
\[
\E e^{i\langle u,Z\rangle}
=\exp\left(i\langle u,m\rangle-\frac12 u^\top\Sigma u\right),
\qquad u\in\R^d.
\]
Prawo tego wektora jest jednoznacznie wyznaczone przez $m$ i $\Sigma$.
\end{stwierdzenie}
\begin{proof}
Z twierdzenia spektralnego mamy $\Sigma=ADA^\top$, gdzie $A$
jest ortogonalna, $D=\operatorname{diag}(\lambda_1,\ldots,\lambda_d)$
i $\lambda_j\ge0$. Niech $G_1,\ldots,G_d$ będą niezależnymi
zmiennymi $\Normal(0,1)$ i połóżmy
$Z=m+AD^{1/2}G$. Wówczas
$\langle u,Z-m\rangle$ jest liniową kombinacją niezależnych
zmiennych normalnych, a zatem jest normalna o wariancji
$u^\top AD A^\top u=u^\top\Sigma u$. Jego funkcję charakterystyczną
otrzymujemy, mnożąc funkcje charakterystyczne składników.
Z jednoznaczności funkcji charakterystycznej wynika jednoznaczność
rozkładu wektora $Z$. Istotnie, dla dowolnego wektora
gaussowskiego o średniej $m$ i kowariancji $\Sigma$ zmienna
$\langle u,Z\rangle$ ma z definicji rozkład normalny,
o średniej $\langle u,m\rangle$ i wariancji $u^\top\Sigma u$.
Jej jednowymiarowa funkcja charakterystyczna w punkcie $1$
daje ten sam wzór, także gdy wariancja jest zerowa.
\end{proof}

\begin{definicja}[Proces gaussowski]
Niech $T\ne\varnothing$ będzie zbiorem indeksów oraz
$(\Omega,\calF,\Prob)$ przestrzenią probabilistyczną. Proces
rzeczywisty $(X_t)_{t\in T}$ na tej przestrzeni nazywamy
\emph{gaussowskim}, jeżeli dla dowolnego $n\ge1$ oraz
$t_1,\ldots,t_n\in T$ wektor $(X_{t_1},\ldots,X_{t_n})$
jest gaussowski. Definiujemy
$m(t):=\E X_t$ i $K(s,t):=\E[(X_s-m(s))(X_t-m(t))]$.
\end{definicja}

\begin{twierdzenie}[Średnia i kowariancja wyznaczają prawo]
Niech $X=(X_t)_{t\in T}$ i $Y=(Y_t)_{t\in T}$ będą rzeczywistymi
procesami gaussowskimi (niekoniecznie na tej samej przestrzeni).
Jeżeli dla wszystkich $s,t\in T$ zachodzą równości
$\E X_t=\E Y_t$ i $\Cov(X_s,X_t)=\Cov(Y_s,Y_t)$, to $X$ i $Y$
mają te same rozkłady skończeniewymiarowe, a więc to samo
prawo na $\R^T$ z $\sigma$-ciałem cylindrycznym.
\end{twierdzenie}
\begin{proof}
Ustalmy $t_1,\ldots,t_n$. Oba odpowiadające im wektory są
gaussowskie, mają wektor średnich
$(m(t_1),\ldots,m(t_n))$ i tę samą macierz kowariancji
$(K(t_i,t_j))_{i,j}$. Z poprzedniego stwierdzenia ich funkcje
charakterystyczne są identyczne. Jednoznaczność funkcji
charakterystycznej daje równość rozkładów. Ponieważ wybór chwil był
dowolny, otrzymujemy równość rozkładów skończeniewymiarowych.
Zdarzenia cylindryczne tworzą $\pi$-układ generujący
$\sigma$-ciało cylindryczne, więc jednoznaczność miar
na tym układzie daje równość praw procesów.
\end{proof}

\begin{uwaga}[Dlaczego warunek wspólnej normalności jest istotny?]
Niech $U\sim\Normal(0,1)$, a $\varepsilon$ będzie od niego
niezależnym symetrycznym znakiem o wartościach $\pm1$.
Wtedy $X=U$, $Y=\varepsilon U$ mają standardowe rozkłady normalne,
a $\Cov(X,Y)=0$, ponieważ $\E\varepsilon=0$.
Jednak $|X|=|Y|$ prawie na pewno. Dla $p=\Prob(|U|\le1)\in(0,1)$
mamy $\Prob(|X|\le1,|Y|\le1)=p\ne p^2$,
zatem zmienne są zależne,
a para $(X,Y)$ nie jest wektorem gaussowskim.
Wektor gaussowski o takiej macierzy kowariancji miałby niezależne
współrzędne, co jest sprzeczne z $|X|=|Y|$.
\end{uwaga}

\section{Jakie funkcje mogą być kowariancjami?}
Zachodzi $K(s,t)=K(t,s)$, ale sama symetria nie wystarcza.
Dla dowolnych $t_1,\ldots,t_n\in T$ i $a_1,\ldots,a_n\in\R$
mamy
\[
\sum_{i,j=1}^n a_i a_jK(t_i,t_j)
=\Var\left(\sum_{i=1}^n a_iX_{t_i}\right)\ge0.
\]
Warunek ten nazywamy \emph{nieujemną określonością jądra} $K$.

\begin{twierdzenie}[Istnienie procesu o zadanej kowariancji]
Niech $T\ne\varnothing$ będzie dowolnym zbiorem,
$m:T\to\R$ dowolną funkcją, zaś $K:T\times T\to\R$
niech spełnia warunki: $K(s,t)=K(t,s)$ oraz
$\sum_{i,j=1}^n a_i a_jK(t_i,t_j)\ge0$ dla wszystkich
$n\in\N$, $t_1,\ldots,t_n\in T$ i $a_1,\ldots,a_n\in\R$.
Wtedy istnieje przestrzeń probabilistyczna i rzeczywisty proces
gaussowski $X=(X_t)_{t\in T}$, dla którego
$\E X_t=m(t)$ oraz $\Cov(X_s,X_t)=K(s,t)$.
\end{twierdzenie}
\begin{proof}
Dla każdej skończonej uporządkowanej rodziny chwil
$\mathbf t=(t_1,\ldots,t_n)$ wybierzmy rozkład gaussowski
$\mu_{\mathbf t}=\Normal_n(m_{\mathbf t},\Sigma_{\mathbf t})$,
gdzie $m_{\mathbf t}=(m(t_1),\ldots,m(t_n))$,
$\Sigma_{\mathbf t}=(K(t_i,t_j))_{i,j}$. Jest to dobrze określony
rozkład, bo macierz kowariancji jest nieujemnie określona.

Trzeba sprawdzić zgodność tej rodziny. Usunięcie części współrzędnych
z wektora gaussowskiego daje wektor gaussowski, którego średnia i
kowariancja są odpowiednimi podwektorami i podmacierzami.
Zmiana kolejności współrzędnych jedynie permutuje średnią i macierz.
Przy powtórzeniu chwili $t_i=t_j$ różnica odpowiednich współrzędnych
ma wariancję $K(t_i,t_i)+K(t_j,t_j)-2K(t_i,t_j)=0$ i średnią zero;
zatem współrzędne są równe prawie na pewno.
Warunki zgodności obowiązują więc także dla powtórzonych chwil.

Możemy teraz zastosować udowodnione w rozdziale~\ref{ch:konstrukcja}
twierdzenie Kołmogorowa o istnieniu procesu, które daje proces
z zadanymi rozkładami skończeniewymiarowymi. Wszystkie te rozkłady
są gaussowskie, a ich średnie i kowariancje są równe $m,K$.
\end{proof}

\begin{stwierdzenie}[Reprezentacja jądra kowariancji w przestrzeni Hilberta]
Dla symetrycznego jądra $K:T\times T\to\R$ nieujemna
określoność jest równoważna istnieniu rzeczywistej
przestrzeni Hilberta $H$ i odwzorowania $\Phi:T\to H$
takiego, że
\[
K(s,t)=\langle\Phi(s),\Phi(t)\rangle_H.
\]
\end{stwierdzenie}
\begin{proof}
Dla jądra nieujemnie określonego poprzednie twierdzenie
daje wycentrowany proces $X$ o kowariancji $K$.
Wystarczy przyjąć $H=L^2(\Omega,\calF,\Prob)$
i $\Phi(t)=X_t$ (klasa równości prawie na pewno).
Wówczas $\langle X_s,X_t\rangle_H=K(s,t)$.
Odwrotnie, dla takiej reprezentacji
\[
\sum_{i,j=1}^n a_ia_jK(t_i,t_j)
=\left\|\sum_{i=1}^n a_i\Phi(t_i)\right\|_H^2\ge0.
\]
Możemy ograniczyć $H$ do domkniętej przestrzeni
liniowej rozpiętej przez $\Phi(T)$.
\end{proof}

\begin{definicja}[Kanoniczna pseudometryka kowariancji]
Dla procesu $X$ o skończonych drugich momentach
definiujemy
\[
d_K(s,t)=\sqrt{K(s,s)+K(t,t)-2K(s,t)}.
\]
Jest to pseudometryka: jest symetryczna, nieujemna,
znika na przekątnej i spełnia nierówność trójkąta,
lecz może znikać dla różnych indeksów.
\end{definicja}
Istotnie, po wycentrowaniu procesu mamy
$d_K(s,t)=\|(X_s-m(s))-(X_t-m(t))\|_{L^2}$,
więc wszystkie własności wynikają z własności normy.
Równość $d_K(s,t)=0$ oznacza równość wycentrowanych
zmiennych prawie na pewno, niekoniecznie równość
ich średnich. Na zbiorze klas równoważności relacji
$s\sim t$ wtedy i tylko wtedy, gdy $d_K(s,t)=0$,
pseudometryka indukuje metrykę: jej niezależność
od reprezentantów wynika z nierówności trójkąta.

\begin{przyklad}[Losowa prosta i losowa fala]
Niech $A,B$ będą niezależne o rozkładzie $\Normal(0,1)$.
Proces $X_t=A+tB$ ma średnią zero i kowariancję
$K_X(s,t)=1+st$. Jego trajektorie są prostymi; proces jest
gaussowski, ponieważ każda skończona rodzina jego wartości
jest liniowym obrazem wektora $(A,B)$.

Dla ustalonej $\omega_0>0$ proces
$Y_t=A\cos(\omega_0t)+B\sin(\omega_0t)$ również jest gaussowski,
a z tożsamości $\cos a\cos b+\sin a\sin b=\cos(a-b)$
wynika $K_Y(s,t)=\cos(\omega_0(t-s))$. Średnia jest zerowa,
a kowariancja nie zmienia się przy wspólnym przesunięciu
argumentów, więc z jednoznaczności praw gaussowskich
wynika stacjonarność w sensie rozkładów skończeniewymiarowych, choć każda trajektoria
jest gładką, okresową funkcją. Stacjonarność nie wymaga więc
ani niezależnych wartości, ani nieregularnych trajektorii.
\end{przyklad}

\begin{przyklad}[Kowariancja wykładnicza: proces typu Ornsteina--Uhlenbecka]
Ustalmy $\sigma>0$ i $\alpha>0$. Zdefiniujmy
$K(s,t)=\sigma^2 e^{-\alpha|t-s|}$ dla $s,t\ge0$.
Rzeczywiście jest to jądro nieujemnie określone, ponieważ
\[
K(s,t)=2\alpha\sigma^2\int_{-\infty}^{\min(s,t)}
 e^{-\alpha(s-u)}e^{-\alpha(t-u)}\,\mathrm d\lambda_1(u).
\]
Dla dowolnych $a_i,t_i$ odpowiadająca mu forma kwadratowa jest
całką z nieujemnego kwadratu:
\[
\sum_{i,j}a_ia_jK(t_i,t_j)=2\alpha\sigma^2\int_{-\infty}^\infty
\left(\sum_i a_i\chi_{\{u\le t_i\}}e^{-\alpha(t_i-u)}\right)^2\mathrm d\lambda_1(u)\ge0.
\]
Poprzednie twierdzenie, zastosowane z $m\equiv0$, daje
wycentrowany proces gaussowski $Y$ o tej kowariancji.
Niezmienniczość średniej i kowariancji względem przesunięć
czasu dowodzi jego stacjonarności. Co więcej,
$\E|Y_t-Y_s|^4=3\,[2\sigma^2(1-e^{-\alpha|t-s|})]^2
\le12\sigma^4\alpha^2|t-s|^2$.
Kryterium ciągłości Kołmogorowa
(twierdzenie~\ref{thm:kol-cont}) zapewnia istnienie ciągłej
modyfikacji na każdym $[0,n]$. Na częściach wspólnych
modyfikacje są nierozróżnialne, co wynika z ich ciągłości
i równości prawie na pewno w każdej wymiernej chwili.
Po przecięciu przeliczalnej rodziny zbiorów pełnej miary
sklejamy je w jedną ciągłą modyfikację na $[0,\infty)$,
a poza tym zbiorem przyjmujemy trajektorię zerową. Nie użyliśmy przy tym całki Itô, którą zdefiniujemy
dopiero w dalszych rozdziałach.
\end{przyklad}

\section{Niezależność i warunkowanie dla procesów gaussowskich}
\begin{twierdzenie}[Nieskorelowane podwektory gaussowskie są niezależne]
Niech $(U,V)$ będzie wektorem gaussowskim o wartościach
w $\R^k\times\R^m$. Jeżeli $\Cov(U_i,V_j)=0$
dla wszystkich $i,j$, to wektory $U$ i $V$ są niezależne.
\end{twierdzenie}
\begin{proof}
Macierz kowariancji wektora $(U,V)$ jest blokowo diagonalna:
$\Sigma=\operatorname{diag}(\Sigma_U,\Sigma_V)$.
Jego funkcja charakterystyczna w punkcie $(u,v)$ jest zatem
iloczynem funkcji charakterystycznych $U$ i $V$,
ponieważ wyrażenie kwadratowe rozdziela się na sumę
$u^\top\Sigma_Uu+v^\top\Sigma_Vv$.
Z jednoznaczności funkcji charakterystycznych wnioskujemy,
że rozkład $(U,V)$ jest iloczynem praw brzegowych.
\end{proof}

\begin{twierdzenie}[Warunkowy rozkład gaussowski]
Niech $(X,Y)$ będzie dwuwymiarowym wektorem gaussowskim,
$\E X=m_X$, $\E Y=m_Y$, $\Var X=\sigma_X^2$,
$\Var Y=\sigma_Y^2>0$ i $\Cov(X,Y)=c$.
Wówczas
\[
X\mid(Y=y)\ \sim\ \Normal\left(
 m_X+\frac{c}{\sigma_Y^2}(y-m_Y),\,
 \sigma_X^2-\frac{c^2}{\sigma_Y^2}\right)
\]
stanowi wersję regularnego rozkładu warunkowego $X$ przy zadanej
wartości $Y=y$. W szczególności
$\E[X\mid\sigma(Y)]=m_X+\frac{c}{\sigma_Y^2}(Y-m_Y)$
i $\Var(X\mid\sigma(Y))=\sigma_X^2-c^2/\sigma_Y^2$ p.n.
\end{twierdzenie}
\begin{proof}
Połóżmy $a=c/\sigma_Y^2$ oraz
$Z=X-m_X-a(Y-m_Y)$. Para $(Z,Y)$ jest gaussowska,
a $\Cov(Z,Y)=c-a\sigma_Y^2=0$.
Z poprzedniego twierdzenia $Z$ i $Y$ są niezależne.
Bezpośredni rachunek daje $\E Z=0$ i
$\Var Z=\sigma_X^2-c^2/\sigma_Y^2\ge0$.
Otrzymaliśmy rozkład $X=m_X+a(Y-m_Y)+Z$,
w którym $Z$ jest niezależna od $Y$.
Dla dowolnej ograniczonej mierzalnej funkcji $f$ i dowolnego
zbioru borelowskiego $A$ niezależność oraz twierdzenie Fubiniego
dają
\[
\E[f(X)\chi_{\{Y\in A\}}]
=\int_A\left(\int_\R f(m_X+a(y-m_Y)+z)\,\mathrm d\Law(Z)(z)\right)
\mathrm d\Law(Y)(y).
\]
Dla $f=\chi_B$ z borelowskim $B$ wewnętrzna całka
jest borelowską funkcją $y$ (z mierzalności funkcji
podcałkowej i twierdzenia Tonellego). Dla ustalonego $y$
jest to prawdopodobieństwo zbioru $B$ względem rozkładu
normalnego. Otrzymujemy zatem jądro probabilistyczne.
Dla ogólnego $f$ wewnętrzna całka jest całką względem
zadeklarowanego rozkładu normalnego (także w przypadku wariancji zero), co dowodzi
własności rozkładu warunkowego. Dla $f(x)=x$ i $x^2$
stosujemy obcięcia i przejście do granicy, wykorzystując
całkowalność drugich momentów, uzyskując obie formuły.
\end{proof}

\begin{uwaga}[Wiele obserwacji i problem najlepszego przewidywania]
Jeżeli $Y\in\R^m$ jest wektorem gaussowskim z odwracalną macierzą
kowariancji $\Sigma_Y$, a $(X,Y)$ jest wspólnie gaussowski,
to ten sam argument z resztą
$Z=X-m_X-\Cov(X,Y)\Sigma_Y^{-1}(Y-m_Y)$ dowodzi wzoru
\[
\E[X\mid\sigma(Y)]
=m_X+\Cov(X,Y)\Sigma_Y^{-1}(Y-m_Y).
\]
Reszta jest nieskorelowana z każdą współrzędną $Y$,
więc niezależna od całego $Y$. To objaśnia, dlaczego
w modelu gaussowskim \emph{najlepsza predykcja średniokwadratowa}
jest funkcją afiniczną zaobserwowanych wartości. Istotnie,
dla $g(Y)\in L^2$ różnica $X-\E[X\mid\sigma(Y)]$
jest ortogonalna do każdej $\sigma(Y)$-mierzalnej
zmiennej z $L^2$, zatem
\[
\E|X-g(Y)|^2=\E|X-\E[X\mid\sigma(Y)]|^2
+\E|\E[X\mid\sigma(Y)]-g(Y)|^2.
\]
Minimum jest osiągane dokładnie przez wersje warunkowej
wartości oczekiwanej.
Gdy $\Sigma_Y$ jest osobliwa, należy uwzględnić liniowe
zależności między obserwacjami. Zamiast odwrotności
zastosujemy poniżej pseudoodwrotność Moore'a--Penrose'a.
\end{uwaga}

\begin{stwierdzenie}[Kryterium Markowa dla procesu gaussowskiego]
Niech $X=(X_t)_{t\in I}$ będzie rzeczywistym, wycentrowanym
procesem gaussowskim na przedziale $I$.
Wówczas $X$ jest procesem Markowa względem swojej
filtracji naturalnej wtedy i tylko wtedy, gdy
\[
K(s,u)=K(s,t)K(t,t)^+K(t,u)
\quad\text{dla wszystkich }s<t<u\text{ w }I,
\]
gdzie $a^+=1/a$ dla $a>0$, a $0^+=0$.
W szczególności przy $K(t,t)=0$ warunek ma postać
$K(s,u)=0$ dla wszystkich $s<t<u$.
\end{stwierdzenie}
\begin{proof}
Ustalmy $t$ i oznaczmy $a_v=K(v,t)K(t,t)^+$.
Jeżeli $K(t,t)=0$, to $X_t=0$ prawie na pewno
i $K(v,t)=0$ dla wszystkich $v$ z nierówności
Cauchy'ego--Schwarza. W obu przypadkach
$a_v K(t,t)=K(v,t)$.
Reszty $R_s=X_s-a_sX_t$ oraz $R_u=X_u-a_uX_t$
tworzą wektor wspólnie gaussowski z $X_t$ i jako para
są niezależne od $X_t$ (obie kowariancje z $X_t$ są zerowe). Obliczenie daje
$\Cov(R_s,R_u)=K(s,u)-K(s,t)K(t,t)^+K(t,u)$.

Jeżeli zachodzi wskazana tożsamość dla każdych $s<t<u$,
to dla dowolnych skończonych zbiorów chwil
$s_1,\ldots,s_k<t<u_1,\ldots,u_m$ trzy bloki
$(R_{s_j})_j$, $(R_{u_l})_l$ oraz $X_t$ mają wspólnie
gaussowski rozkład o blokowo diagonalnej macierzy
kowariancji. Funkcja charakterystyczna rozkłada się
na iloczyn trzech czynników, więc bloki są łącznie
niezależne. W szczególności wektor przyszłych reszt
jest niezależny od $\sigma(X_t,R_s:s<t)$: dla
zdarzeń cylindrycznych wynika to z powyższej
niezależności, a argument $\pi$--$\lambda$ rozszerza
ją na całe $\sigma$-ciało. Ponieważ
$\sigma(X_t,R_s:s<t)=\sigma(X_s:s\le t)$,
dla każdej ograniczonej borelowskiej funkcji $h$ mamy
\[
\E[h(X_u)\mid\sigma(X_s:s\le t)]
=H(X_t),\quad H(x)=\E h\left(
a_ux+R_u\right).
\]
Mierzalność $H$ wynika z twierdzenia Fubiniego,
a tożsamość z niezależności $R_u$ od przeszłości.
Warunkując jeszcze względem $\sigma(X_t)$,
otrzymujemy własność Markowa.

Odwrotnie, własność Markowa daje warunkową niezależność
$X_s$ i $X_u$ względem $X_t$ dla $s<t<u$: dla
ograniczonych funkcji przeszłości i przyszłości stosujemy
własność wieży i warunkowanie przyszłości względem
$\sigma(X_r:r\le t)$; funkcje nieograniczone dopuszczamy
następnie przez obcięcie w $L^2$. W szczególności $\E[X_u\mid\calF_t^X]
=\E[X_u\mid\sigma(X_t)]=a_uX_t$,
gdzie $\calF_t^X=\sigma(X_s:s\le t)$, a przejście
od funkcji ograniczonych do $X_u$ jest dozwolone przez
zbieżność w $L^2$ i kontraktywność warunkowej wartości
oczekiwanej. Ostatnia równość wynika z niezależności
wycentrowanej reszty $R_u$ od $X_t$, także gdy $X_t$
jest stała prawie na pewno. Zatem $\E[R_u\mid\calF_t^X]=0$.
Ponieważ $R_s$ jest $\calF_t^X$-mierzalna i należy do $L^2$,
przez obcięcie $R_s$ i nierówność Cauchy'ego--Schwarza
otrzymujemy $\E[R_sR_u]=0$. Reszty są wycentrowane,
więc $\Cov(R_s,R_u)=0$, czyli żądana tożsamość.
\end{proof}


\begin{twierdzenie}[Warunkowanie wielowymiarowe: pełny wzór]
Niech $(X,Y)$ będzie wspólnie gaussowskim wektorem,
$X\in\R^k$, $Y\in\R^m$, o średnich $m_X,m_Y$ i blokowej
macierzy kowariancji
$\left(\begin{smallmatrix}\Sigma_X&C\\ C^\top&\Sigma_Y\end{smallmatrix}\right)$.
Niech $\Sigma_Y^+$ oznacza pseudoodwrotność
Moore'a--Penrose'a macierzy kowariancji $\Sigma_Y$:
jeżeli $\Sigma_Y=U\operatorname{diag}(d_1,\ldots,d_m)U^\top$
jest rozkładem spektralnym, to
\[
\Sigma_Y^+=U\operatorname{diag}(d_1^+,\ldots,d_m^+)U^\top,
\qquad
d_j^+=\begin{cases}1/d_j,&d_j>0,\\0,&d_j=0.\end{cases}
\]
Definicja nie zależy od wyboru bazy własnej, ponieważ
operator działa skalarnie na każdej przestrzeni własnej.
Wtedy regularny rozkład warunkowy $X$ przy $Y=y$ ma
wersję normalną z parametrami
\[
\mu(y)=m_X+C\Sigma_Y^+(y-m_Y),\qquad
V=\Sigma_X-C\Sigma_Y^+C^\top.
\]
Macierz $V$ jest nieujemnie określona, także wtedy,
gdy rozkład warunkowy jest zdegenerowany.
Jeżeli $\Sigma_Y$ jest odwracalna, to $\Sigma_Y^+=\Sigma_Y^{-1}$.
\end{twierdzenie}
\begin{proof}
Dla $v\in\ker\Sigma_Y$ zmienna $v^\top(Y-m_Y)$
ma średnią i wariancję zero, więc jest zerowa prawie
na pewno. W konsekwencji
$Cv=\E[(X-m_X)v^\top(Y-m_Y)]=0$.
Zatem $\ker\Sigma_Y\subseteq\ker C$.
Ze wzoru spektralnego wynika, że
$\Sigma_Y^+\Sigma_Y=\Sigma_Y\Sigma_Y^+$ jest rzutem
ortogonalnym na obraz $\Sigma_Y$, a
$\Sigma_Y^+\Sigma_Y\Sigma_Y^+=\Sigma_Y^+$.
Ponieważ $C$ znika na jądrze $\Sigma_Y$,
mamy $C\Sigma_Y^+\Sigma_Y=C$.

Połóżmy $A=C\Sigma_Y^+$ i
$R=X-m_X-A(Y-m_Y)$. Wektor $(R,Y)$ jest gaussowski,
bo jest afinicznym obrazem wektora $(X,Y)$.
Ponieważ
$\Cov(R,Y)=C-A\Sigma_Y=0$,
wektory $R,Y$ są niezależne. Ponadto $\E R=0$.
Rozwijając macierz kowariancji, otrzymujemy
\[
\Cov(R,R)=\Sigma_X-AC^\top-CA^\top+A\Sigma_YA^\top
=\Sigma_X-C\Sigma_Y^+C^\top=V.
\]
Jako macierz kowariancji $V$ musi być nieujemnie określona.
Ponadto $X=m_X+A(Y-m_Y)+R$. Dla ograniczonej funkcji
borelowskiej $h$ na $\R^k$ i zbioru borelowskiego
$D\subset\R^m$, z niezależności oraz twierdzenia Fubiniego,
\[
\E[h(X)\chi_{\{Y\in D\}}]
=\int_D\int_{\R^k}h(m_X+A(y-m_Y)+r)
\,\mathrm d\Law(R)(r)\,\mathrm d\Law(Y)(y).
\]
Dla funkcji charakterystycznych zbiorów borelowskich
mierzalność wewnętrznej całki względem $y$ wynika
z twierdzenia Tonellego; dla ustalonego $y$ jest ona
prawdopodobieństwem względem prawa przesuniętego $R$.
To dokładnie własność regularnego rozkładu warunkowego
o podanych parametrach. Rozkład $Y$ jest skoncentrowany
na $m_Y+\operatorname{im}\Sigma_Y$: wystarczy wybrać
skończoną bazę jądra i przeciąć zbiory pełnej miary,
na których odpowiednie iloczyny skalarne z $Y-m_Y$
są zerowe. Wartości jądra poza tą podprzestrzenią
nie są wyznaczone przez warunek rozkładu warunkowego;
podany wzór wybiera jedną wersję na całym $\R^m$.
\end{proof}

\begin{przyklad}[Powtórzona obserwacja]
Niech $U,V$ będą niezależnymi zmiennymi standardowo
normalnymi, $X=U+V$, a $Y=(U,U)^\top$.
Wtedy
\[
\Sigma_Y=\begin{pmatrix}1&1\\1&1\end{pmatrix},\qquad
\Sigma_Y^+=\tfrac14\begin{pmatrix}1&1\\1&1\end{pmatrix},
\qquad C=\begin{pmatrix}1&1\end{pmatrix}.
\]
Wzór daje $X\mid(Y=y)\sim
\Normal((y_1+y_2)/2,1)$. Na prostej $y_1=y_2$
jest to rozkład $\Normal(y_1,1)$, zgodnie z niezależnością
$V$ od $U$. Powtórzenie tej samej obserwacji nie
zmniejsza wariancji błędu predykcji.
\end{przyklad}

\begin{twierdzenie}[Predykcja względem dowolnej rodziny obserwacji gaussowskich]
\label{thm:gauss-prediction}
Niech $X$ oraz $(Y_j)_{j\in J}$ będą wspólnie gaussowską
rodziną rzeczywistych zmiennych losowych, to znaczy
każda skończona podrodzina jest wektorem gaussowskim.
Niech $\mathcal G=\sigma(Y_j:j\in J)$ i niech
\[
H=\overline{\operatorname{span}\{Y_j-\E Y_j:j\in J\}}^{L^2}
\]
będzie domkniętą podprzestrzenią rzeczywistej przestrzeni
Hilberta $L^2(\Omega,\calF,\Prob)$. Oznaczmy przez $P_H$
rzut ortogonalny na $H$. Wtedy
\[
\E[X\mid\mathcal G]=\E X+P_H(X-\E X).
\]
Reszta $R=X-\E X-P_H(X-\E X)$ jest wycentrowana,
gaussowska i niezależna od $\mathcal G$. W szczególności
wariancja warunkowa jest stała i równa $\E R^2$.
\end{twierdzenie}
\begin{proof}
Niech $Z=P_H(X-\E X)$. Z definicji domknięcia
istnieje ciąg skończonych kombinacji liniowych
wycentrowanych obserwacji $Z_n\to Z$ w $L^2$.
Z podciągu zbieżnego prawie na pewno otrzymujemy
$\mathcal G$-mierzalną wersję $Z$ (definiujemy
granicę na mierzalnym zbiorze jej istnienia, a poza nim
przyjmujemy zero). Ponieważ $\E Z_n=0$, ciągłość
wartości oczekiwanej w $L^2$ daje $\E Z=0$.

Ustalmy skończenie wiele indeksów $j_1,\ldots,j_m$.
Każdy wektor $(X-\E X-Z_n,Y_{j_1},\ldots,Y_{j_m})$
jest gaussowski. Dowolna jego kombinacja liniowa
$S_n$ zbiega w $L^2$ do odpowiadającej kombinacji $S$
wektora $(R,Y_{j_1},\ldots,Y_{j_m})$.
Wtedy $\E S_n\to\E S$ oraz $\Var S_n\to\Var S$,
a dla $u\in\R$ nierówność
$|e^{iuS_n}-e^{iuS}|\le |u||S_n-S|$
daje zbieżność funkcji charakterystycznych.
W granicy otrzymujemy funkcję charakterystyczną
$\Normal(\E S,\Var S)$, także przy wariancji zero.
Z jednoznaczności funkcji charakterystycznej $S$
jest gaussowska. Tym samym $(R,Y_{j_1},\ldots,Y_{j_m})$
jest wspólnie gaussowski.

Z własności rzutu ortogonalnego mamy
$\E[R(Y_j-\E Y_j)]=0$ dla wszystkich $j\in J$.
Nieskorelowanie i wspólna gaussowskość dają niezależność
$R$ od każdej skończonej rodziny obserwacji. Zdarzenia
cylindryczne tworzą $\pi$-układ generujący $\mathcal G$.
Dla ustalonego borelowskiego zbioru $B$ stosujemy
argument $\pi$--$\lambda$ do miar
$A\mapsto\Prob(A\cap\{R\in B\})$ oraz
$A\mapsto\Prob(A)\Prob(R\in B)$ na $\mathcal G$.
Otrzymujemy niezależność $R$ od $\mathcal G$,
a więc $\E[R\mid\mathcal G]=0$ oraz
$\E[R^2\mid\mathcal G]=\E R^2$.
Z rozkładu $X=\E X+Z+R$ wynika teza.
\end{proof}

\begin{uwaga}[Interpretacja geometryczna]
Predykcja gaussowska jest rzutem ortogonalnym nawet wtedy,
gdy rodzina obserwacji jest nieprzeliczalna. Konkretny rzut
można przybliżać kombinacjami skończenie wielu obserwacji.
Nie jest potrzebne założenie o przeliczalności całej
rodziny. W przypadku skończonym współczynniki rzutu
spełniają równania normalne $A\Sigma_Y=C$;
rozwiązanie $A=C\Sigma_Y^+$ prowadzi do poprzedniego
wzoru. Niezależność reszty od wszystkich obserwacji
jest mocniejsza niż sama ortogonalność w $L^2$.
\end{uwaga}

\begin{przyklad}[Prawo przejścia procesu Ornsteina--Uhlenbecka]
Dla wycentrowanego procesu $Y$ o kowariancji
$K(s,t)=\sigma^2e^{-\alpha|t-s|}$, przy $s<t<u$,
\[
K(s,t)K(t,t)^{-1}K(t,u)
=\sigma^2e^{-\alpha(t-s)}e^{-\alpha(u-t)}
=K(s,u).
\]
Kryterium Markowa daje więc własność Markowa względem
filtracji naturalnej. Dla $h\ge0$ rozkład warunkowy
$Y_{t+h}$ względem całej przeszłości do $t$ ma postać
\[
\Normal\bigl(e^{-\alpha h}Y_t,\,
\sigma^2(1-e^{-2\alpha h})\bigr).
\]
Wynika to z warunkowego rozkładu pary $(Y_{t+h},Y_t)$.
Przy $h=0$ otrzymujemy rozkład skupiony w $Y_t$.
Dla $h\to\infty$ średnia warunkowa dąży do zera,
a wariancja do $\sigma^2$. Warunkowy rozkład przy
$Y_t=x$ dąży słabo do $\Normal(0,\sigma^2)$
dla każdego $x$, co można zobaczyć, przedstawiając go
jako prawo $e^{-\alpha h}x+\sigma\sqrt{1-e^{-2\alpha h}}G$,
gdzie $G\sim\Normal(0,1)$, i stosując twierdzenie
o zbieżności zmajoryzowanej do ograniczonych funkcji
ciągłych. Jest to przykład zaniku pamięci o stanie
początkowym w modelu stacjonarnym.
\end{przyklad}

\section{Od kowariancji do ciągłości trajektorii}

\begin{stwierdzenie}[Ciągłość średniokwadratowa a kowariancja]
Niech $I$ będzie przedziałem i niech $X=(X_t)_{t\in I}$
będzie procesem o skończonych drugich momentach.
Następujące warunki są równoważne:
\begin{itemize}
\item odwzorowanie $t\mapsto X_t$ jest ciągłe w $L^2$;
\item funkcja średniej $m$ i jądro kowariancji $K$
są ciągłe (odpowiednio na $I$ i $I\times I$).
\end{itemize}
\end{stwierdzenie}
\begin{proof}
Rozkład przyrostu na część wycentrowaną i średnią daje
\[
\E|X_t-X_s|^2=|m(t)-m(s)|^2+d_K(s,t)^2.
\]
Stąd ciągłość $m,K$ implikuje ciągłość w $L^2$.
Odwrotnie, z nierówności Cauchy'ego--Schwarza
$|\E X_t-\E X_s|\le\|X_t-X_s\|_{L^2}$.
Ciągły jest więc również proces wycentrowany
$Z_t=X_t-m(t)$ jako odwzorowanie do $L^2$.
Jeżeli $(s_n,t_n)\to(s,t)$, to
\[
|K(s_n,t_n)-K(s,t)|
\le\|Z_{s_n}-Z_s\|_{L^2}\|Z_{t_n}\|_{L^2}
+\|Z_s\|_{L^2}\|Z_{t_n}-Z_t\|_{L^2}.
\]
Normy $\|Z_{t_n}\|_{L^2}$ są ograniczone, gdyż
$Z_{t_n}\to Z_t$ w $L^2$. Prawa strona dąży do zera.
\end{proof}

\begin{wniosek}[Ciągłość stochastyczna procesu gaussowskiego]
Dla rzeczywistego procesu gaussowskiego na przedziale
ciągłość stochastyczna jest równoważna ciągłości
średniokwadratowej.
\end{wniosek}
\begin{proof}
Ciągłość w $L^2$ implikuje ciągłość stochastyczną
na mocy nierówności Markowa. Odwrotnie, niech
$t_n\to t$ i $X_{t_n}-X_t\to0$ według prawdopodobieństwa.
Oznaczmy średnią i wariancję przyrostu przez $a_n,b_n$.
Zbieżność według prawdopodobieństwa daje zbieżność
jego funkcji charakterystycznej do $1$ w każdym
punkcie $u\in\R$ (można zastosować ograniczoność
funkcji wykładniczej i kryterium podciągów).
Dla $u=1$ moduł tej funkcji wynosi $e^{-b_n/2}$,
więc $b_n\to0$. Wycentrowane przyrosty zbiegają
zatem do zera w $L^2$, a więc również według
prawdopodobieństwa. Odejmując je od przyrostów,
otrzymujemy $a_n\to0$ według prawdopodobieństwa;
ponieważ $a_n$ są deterministyczne, jest to
zwykła zbieżność liczb. W konsekwencji
$\E|X_{t_n}-X_t|^2=a_n^2+b_n\to0$.
\end{proof}

\begin{twierdzenie}[Różniczkowanie średniokwadratowe]
\label{thm:gauss-ms-derivative}
Niech $X$ będzie procesem o skończonych drugich
momentach na otwartym przedziale $I$. Załóżmy, że
$m\in C^1(I)$ oraz $K\in C^2(I\times I)$.
Wówczas dla każdego $t\in I$ istnieje pochodna
średniokwadratowa
\[
D_t=\lim_{h\to0}\frac{X_{t+h}-X_t}{h}
\quad\text{w }L^2.
\]
Zachodzą wzory
\[
\E D_t=m'(t),\qquad
\Cov(D_s,D_t)=\partial_s\partial_t K(s,t),
\qquad
\Cov(X_s,D_t)=\partial_t K(s,t).
\]
Proces $D$ jest ciągły średniokwadratowo.
Jeżeli $X$ jest gaussowski, to rodzina jego wartości
i pochodnych średniokwadratowych jest wspólnie
gaussowska.
\end{twierdzenie}
\begin{proof}
Niech $Z_t=X_t-m(t)$ i
$A_h(t)=(Z_{t+h}-Z_t)/h$ dla $h\ne0$, $t+h\in I$.
Dwukrotne zastosowanie podstawowego twierdzenia
rachunku całkowego daje dla dopuszczalnych $h,k$
\[
\E[A_h(s)A_k(t)]
=\int_0^1\int_0^1
\partial_s\partial_t K(s+uh,t+vk)
\,\mathrm d\lambda_1(u)\,\mathrm d\lambda_1(v).
\]
Wzór obowiązuje także dla ujemnych $h,k$, gdyż
stosujemy parametryzację odcinków przez $[0,1]$.
Z ciągłości pochodnej mieszanej prawa strona dąży
do $\partial_s\partial_t K(s,t)$ przy $h,k\to0$.
W szczególności dla $s=t$, rozwijając
$\E|A_h(t)-A_k(t)|^2$, otrzymujemy granicę zero.
Z zupełności $L^2$ wynika istnienie granicy $A(t)$
w $L^2$ (najpierw dla dowolnego ciągu $h_n\to0$,
a warunek Cauchy'ego zapewnia tę samą granicę dla
wszystkich $h\to0$). Mamy $\E A(t)=0$.
Dodając iloraz różnicowy $m$, otrzymujemy
$D_t=A(t)+m'(t)$. Granica iloczynów skalarnych
w $L^2$ daje wzór na kowariancję pochodnych.
Ponadto z ciągłości iloczynu skalarnego
\[
\Cov(X_s,D_t)=\lim_{h\to0}\frac{K(s,t+h)-K(s,t)}{h}
=\partial_t K(s,t).
\]
Funkcja średniej procesu $D$ oraz jego jądro kowariancji
są ciągłe na mocy założeń, więc poprzednie stwierdzenie
daje ciągłość średniokwadratową $D$.
Jeżeli $X$ jest gaussowski, każda skończona kombinacja
liniowa wartości i ilorazów różnicowych jest normalna.
Przechodząc do granicy w $L^2$, stosujemy argument
z funkcjami charakterystycznymi użyty w dowodzie
twierdzenia~\ref{thm:gauss-prediction}. Otrzymujemy wspólną gaussowskość.
\end{proof}

\begin{przyklad}[Gładkość średniokwadratowa losowej fali i procesu Ornsteina--Uhlenbecka]
Dla losowej fali o kowariancji
$K(s,t)=\cos(\omega_0(t-s))$ mamy
$\partial_s\partial_t K(s,t)=\omega_0^2\cos(\omega_0(t-s))$.
Pochodna średniokwadratowa jest tu także pochodną
trajektorii, gdyż bezpośrednio różniczkujemy
$A\cos(\omega_0t)+B\sin(\omega_0t)$, a zbieżność
współczynników daje zbieżność w $L^2$.
Dla procesu Ornsteina--Uhlenbecka natomiast
\[
\E\left|\frac{Y_{t+h}-Y_t}{h}\right|^2
=\frac{2\sigma^2(1-e^{-\alpha|h|})}{h^2}
\longrightarrow\infty\qquad(h\to0).
\]
Wynika to z granicy $(1-e^{-x})/x\to1$ przy
$x\downarrow0$. Ilorazy różnicowe nie mogą więc
zbiegać w $L^2$, choć proces ma ciągłą modyfikację.
Brak pochodnej w $L^2$ sam w sobie, dla dowolnego
procesu, nie jest dowodem braku różniczkowalności
trajektorii. Tutaj można jednak powiedzieć więcej:
ilorazy różnicowe procesu Ornsteina--Uhlenbecka
są wycentrowanymi zmiennymi normalnymi o wariancjach
dążących do nieskończoności. Dla każdego $M>0$
prawdopodobieństwo, że ich moduł nie przekracza $M$,
dąży do zera, ponieważ po standaryzacji otrzymujemy
przedziały o długościach dążących do zera, a gęstość
standardowego rozkładu normalnego jest ograniczona.
Nie mogą więc zbiegać według prawdopodobieństwa do
skończonej zmiennej losowej. Co więcej, dla dowolnego
deterministycznego ciągu $h_n\to0$, $h_n\ne0$,
zdarzenie zbieżności ilorazów do skończonej granicy
ma prawdopodobieństwo zero: zawiera się w sumie po
$M\in\N$ zdarzeń
$\bigcup_N\bigcap_{n\ge N}\{|(Y_{t+h_n}-Y_t)/h_n|\le M\}$,
z których każde ma miarę zero.
Zatem dla każdego ustalonego $t>0$ prawdopodobieństwo
istnienia pochodnej trajektorii w $t$ wynosi zero.
Nie jest to jeszcze twierdzenie o jednoczesnym braku
pochodnej we wszystkich chwilach.
\end{przyklad}

\begin{lemat}[Moment przyrostu gaussowskiego]
Jeżeli $X$ jest rzeczywistym procesem gaussowskim,
a $X_t-X_s$ ma średnią zero oraz wariancję $v(s,t)$,
to dla każdego $p>0$
\[
\E|X_t-X_s|^p=c_p\,v(s,t)^{p/2},\qquad
c_p=\frac{2^{p/2}\Gamma((p+1)/2)}{\sqrt\pi}.
\]
\end{lemat}
\begin{proof}
Gdy $v(s,t)=0$, przyrost wynosi zero prawie na pewno.
W przeciwnym razie $(X_t-X_s)/\sqrt{v(s,t)}$ ma standardowy
rozkład normalny. Podstawienie w całce względem jego gęstości
i zastosowanie definicji funkcji gamma daje
$\E|G|^p=(2/\sqrt{2\pi})\int_0^\infty
 x^p e^{-x^2/2}\mathrm d\lambda_1(x)=c_p$.
\end{proof}

\begin{wniosek}[Praktyczne kryterium ciągłości procesu gaussowskiego]
Niech $X=(X_t)_{t\in[0,T]}$, $T>0$, będzie procesem gaussowskim.
Załóżmy, że istnieją $C>0$, $\eta>0$ takie, że
$\E|X_t-X_s|^2\le C|t-s|^{\eta}$ dla wszystkich $s,t\in[0,T]$.
Nie wymagamy stałej funkcji średniej.
Wtedy $X$ ma ciągłą modyfikację. Co więcej, ma modyfikację,
której trajektorie są h\"olderowskie na $[0,T]$
każdego rzędu $0<\gamma<\min(\eta/2,1)$
na jednym zbiorze pełnej miary.
\end{wniosek}
\begin{proof}
Niech $m(t)=\E X_t$. Nierówność Cauchy'ego--Schwarza
daje $|m(t)-m(s)|\le C^{1/2}|t-s|^{\eta/2}$.
Wariancja przyrostu jest nie większa od jego drugiego
momentu. Stosując poprzedni lemat do wycentrowanego
przyrostu oraz nierówność
$|a+b|^p\le2^{p-1}(|a|^p+|b|^p)$, dla $p\ge2$ dostajemy
\[
\E|X_t-X_s|^p\le
2^{p-1}(1+c_p)C^{p/2}|t-s|^{\eta p/2}.
\]
Wybierzmy $p$ tak duże, by $\eta p/2>1$.
Twierdzenie~\ref{thm:kol-cont} daje ciągłą modyfikację,
a także wykładniki H\"oldera mniejsze niż
$(\eta p/2-1)/p=\eta/2-1/p$ (oraz mniejsze niż $1$).
Dążąc z $p$ do nieskończoności i korzystając z
nierozróżnialności ciągłych modyfikacji, otrzymujemy
jedną modyfikację o wszystkich wskazanych wykładnikach
(po ograniczeniu do przeliczalnego ciągu wykładników
zbiegającego do $\min(\eta/2,1)$).
\end{proof}

\begin{uwaga}[Ciągłość kowariancji to nie całe rozwiązanie]
W szczególności, jeżeli proces gaussowski jest wycentrowany,
to $\E|X_t-X_s|^2=K(t,t)+K(s,s)-2K(s,t)$.
Oszacowanie tej różnicy daje konkretne kryterium regularności.
Nie należy jednak wnioskować automatycznie o ciągłości
trajektorii wyłącznie z samej ciągłości funkcji $K$:
kryterium momentowe wymaga kontroli \emph{szybkości}
zanikania wariancji przyrostów. Bardziej ogólne kryteria
wykorzystują między innymi metryczną entropię.
\end{uwaga}

\begin{przyklad}[Ciągła kowariancja bez ciągłej modyfikacji]
Niech $(G_n)_{n\ge1}$ będą niezależnymi zmiennymi
$\Normal(0,1)$. Na każdym przedziale
$I_n=[2^{-(n+1)},2^{-n}]$ niech $\varphi_n$ będzie
ciągłą funkcją łamaną, równą zero na końcach,
równą $1$ w środku $c_n$ i zerową poza $I_n$.
Na $[0,1]$ definiujemy
\[
X_t=\sum_{n\ge1}\frac{G_n\varphi_n(t)}{\sqrt{\log(n+1)}}.
\]
Dla każdego $t$ co najwyżej jeden składnik jest
niezerowy, więc każda skończona rodzina wartości
jest liniowym obrazem skończonego wektora niezależnych
zmiennych normalnych. Proces jest wycentrowany, a
\[
K(s,t)=\sum_{n\ge1}\frac{\varphi_n(s)\varphi_n(t)}{\log(n+1)}.
\]
Szereg jest jednostajnie zbieżny: dla każdego $(s,t)$
co najwyżej jeden składnik jest niezerowy; ponieważ
$0\le\varphi_n\le1$, supremum ogona po $n>N$
nie przekracza $1/\log(N+2)$.
Zatem $K$ jest ciągłe na $[0,1]^2$ i proces jest
ciągły średniokwadratowo.

Pokażemy, że nie istnieje ciągła modyfikacja.
Niech $a_n=\sqrt{\log(n+1)}$. Dla standardowej
zmiennej normalnej $G$ i $a\ge1$ mamy
\[
\Prob(G>a)\ge
\frac1{\sqrt{2\pi}}\int_a^{a+1/a}e^{-x^2/2}\,
\mathrm d\lambda_1(x)
\ge \frac{e^{-3/2}}{\sqrt{2\pi}a}e^{-a^2/2}.
\]
Ostatnia nierówność wynika z
$(a+1/a)^2/2\le a^2/2+3/2$.
Stąd szereg $\sum_n\Prob(G_n>a_n)$ jest rozbieżny,
gdyż jego wyrazy dla dużych $n$ są ograniczone od dołu
przez dodatnią stałą razy
$1/\sqrt{(n+1)\log(n+1)}$, a $\log(n+1)\le n+1$.
Z niezależności zdarzeń i drugiego lematu
Borela--Cantellego otrzymujemy $X_{c_n}>1$
dla nieskończenie wielu $n$ prawie na pewno.
Ponieważ $c_n\to0$ oraz $X_0=0$, ciągłość w zerze
nie zachodzi. Każda modyfikacja zgadza się z $X$
w chwilach $0,c_1,c_2,\ldots$ na jednym zbiorze
pełnej miary, więc również nie może być ciągła
w zerze prawie na pewno.
\end{przyklad}

\section{Most gaussowski: obserwacja końca zmienia cały proces}
W rozdziale~\ref{ch:konstrukcja} skonstruowaliśmy proces
Wienera $W$; w następnym rozdziale zbadamy jego własności.
Przypomnijmy jego charakterystyczną kowariancję
$\Cov(W_s,W_t)=\min(s,t)$ dla $s,t\ge0$.
Ustalmy $T>0$ i zdefiniujmy na $[0,T]$
\[
\beta_t:=W_t-\frac{t}{T}W_T.
\]
Wówczas $\beta_0=\beta_T=0$, a zatem każda trajektoria
\emph{mostu Browna} zaczyna i kończy się w zerze.

\begin{twierdzenie}[Most Browna i warunkowanie na punkt końcowy]
Niech $W$ będzie standardowym procesem Wienera i $T>0$.
Proces $\beta=(\beta_t)_{0\le t\le T}$ jest wycentrowanym
procesem gaussowskim, niezależnym od zmiennej $W_T$,
z funkcją kowariancji
\[
K_\beta(s,t)=\min(s,t)-\frac{st}{T},
\qquad 0\le s,t\le T.
\]
Ponadto prawo procesu $t\mapsto (t/T)y+\beta_t$
na $C([0,T],\R)$ z borelowskim $\sigma$-ciałem
normy supremum, jako funkcja $y$, stanowi wersję
regularnego rozkładu warunkowego całej trajektorii $W$
przy warunku $W_T=y$.
\end{twierdzenie}
\begin{proof}
Każdy skończony wektor wartości $\beta$ jest liniowym obrazem
wektora wartości $W$, więc jest gaussowski.
Ponieważ $\Cov(W_t,W_T)=t$ oraz $\Var W_T=T$,
otrzymujemy $\Cov(\beta_t,W_T)=t-t=0$.
Wspólna gaussowskość wszystkich skończonych wektorów daje
niezależność $(\beta_{t_1},\ldots,\beta_{t_n})$ od $W_T$.
Przechodząc ze zdarzeń cylindrycznych do $\sigma$-ciała generowanego
przez proces, otrzymujemy niezależność całego $\beta$ od $W_T$.
Rozwijając kowariancję, otrzymujemy
\[
\Cov(\beta_s,\beta_t)=\min(s,t)
-\frac{t}{T}s-\frac{s}{T}t+\frac{st}{T^2}T
=\min(s,t)-\frac{st}{T}.
\]
Na koniec równość $W_t=(t/T)W_T+\beta_t$ i niezależność
pozwalają rozwarunkować względem $W_T$ dokładnie tak,
jak w dowodzie wzoru na warunkowy rozkład gaussowski.
Dla każdej ograniczonej borelowskiej funkcji $F$ skończonej
liczby współrzędnych oraz zbioru $A$ zachodzi
\[
\E\bigl[F(W_{t_1},\ldots,W_{t_n})\chi_{\{W_T\in A\}}\bigr]
=\int_A \E\bigl[F((t_1/T)y+\beta_{t_1},\ldots,
(t_n/T)y+\beta_{t_n})\bigr]\,\mathrm d\Law(W_T)(y),
\]
co dowodzi tezy dla cylindrów. Przyjmujemy ciągłe
wersje procesów, z $W_0=0$ punktowo.
Borelowskie $\sigma$-ciało $C([0,T],\R)$ jest generowane
przez ewaluacje w przeliczalnym gęstym zbiorze chwil:
normę supremum różnicy funkcji ciągłych można obliczać
po tym zbiorze, więc każda otwarta kula jest mierzalna
względem ewaluacji. Przestrzeń $C([0,T],\R)$ jest
ośrodkowa (na przykład łamane na siatkach
$kT/n$ o wymiernych wartościach w węzłach tworzą
przeliczalny zbiór gęsty), zatem każdy zbiór otwarty
jest przeliczalną sumą takich kul. Odwrotna inkluzja
wynika z ciągłości ewaluacji. W szczególności $W$
i $\beta$ są mierzalnymi zmiennymi losowymi
o wartościach w tej przestrzeni.

Oznaczmy przez $\nu$ prawo $\beta$ na tej przestrzeni.
Odwzorowanie
$(y,f)\mapsto f+(\cdot/T)y$ jest ciągłe względem
normy supremum. Dla borelowskiego $B\subset C([0,T],\R)$
definiujemy
\[
Q(y,B)=\int\chi_B(f+(\cdot/T)y)\,\mathrm d\nu(f).
\]
Jest to jądro probabilistyczne: mierzalność względem $y$
wynika z twierdzenia Tonellego, a dla ustalonego $y$
jest to miara będąca obrazem $\nu$ przez przesunięcie.
Dla borelowskiego $A\subseteq\R$ równość
\[
\Prob(W\in B,\ W_T\in A)
=\int_A Q(y,B)\,\mathrm d\Law(W_T)(y)
\]
zachodzi dla cylindrów na mocy wcześniejszego rachunku.
Dla ustalonego $A$ obie strony są miarami skończonymi
względem $B$, więc argument $\pi$--$\lambda$ rozszerza
równość na całe borelowskie $\sigma$-ciało.
To dowodzi tezy dla całych trajektorii. Ponieważ
$\Prob(W_T=y)=0$ dla każdego $y$, chodzi o wersję
jądra, a nie o iloraz prawdopodobieństw zdarzeń.
\end{proof}

\begin{center}
\begin{tikzpicture}[x=0.9cm,y=0.8cm,>=Latex]
\draw[->] (-0.3,0)--(10.3,0) node[right] {$t$};
\draw[->] (0,-1.4)--(0,2.0) node[above] {$\beta_t$};
\draw[dashed,gray] (0,0)--(10,0);
\draw[very thick,blue!70!black]
 plot coordinates {(0,0)(0.6,0.5)(1.2,0.2)(1.8,1.1)(2.5,0.35)(3,0.8)(3.7,0.2)(4.5,-0.8)(5.1,-0.35)(5.9,0.9)(6.6,0.2)(7.3,0.85)(8.1,0.25)(8.8,-0.35)(9.4,0.2)(10,0)};
\node[below] at (0,0) {$0$};
\node[below] at (10,0) {$T$};
\node[above right,blue!70!black] at (2,1.15) {przykładowy most};
\end{tikzpicture}
\par\smallskip\small Schematyczna łamana ilustrująca most Browna: oba końce są
przypięte do zera. Łamana jest jedynie ilustracją, nie
rzeczywistą trajektorią procesu.
\end{center}

\begin{zadanie}[Predykcja wartości procesu Wienera na podstawie dwóch obserwacji]
Niech $W$ będzie procesem Wienera, $0<s<t$ oraz $s\le u\le t$.
Obliczyć $\E[W_u\mid\sigma(W_s,W_t)]$ oraz wariancję warunkową.

\emph{Rozwiązanie.} Z niezależności przyrostów $W_s$, $W_u-W_s$ i $W_t-W_u$
są łącznie niezależne, z uwzględnieniem zerowych
przyrostów na końcach przedziału. Zatem pierwsza
współrzędna wektora $(W_s,W_u-W_s,W_t-W_s)$
jest niezależna od pary pozostałych. W szczególności
$\E[W_u\mid W_s,W_t]=W_s+\E[W_u-W_s\mid W_t-W_s]$.
Obserwowane dwa przyrosty są wspólnie gaussowskie,
a $\Cov(W_u-W_s,W_t-W_s)=u-s$ i
$\Var(W_t-W_s)=t-s$.
Wzór na warunkowy rozkład gaussowski daje
\[
\E[W_u\mid W_s,W_t]
=W_s+\frac{u-s}{t-s}(W_t-W_s),\qquad
\Var(W_u\mid W_s,W_t)=\frac{(u-s)(t-u)}{t-s}.
\]
Zatem predykcją jest interpolacja liniowa między pomiarami,
a wariancja błędu jest zerowa na końcach przedziału
i osiąga maksimum $(t-s)/4$ w jego środku.
\end{zadanie}

\section{Podsumowanie: dlaczego procesy gaussowskie przygotowują ruch Browna?}
Mamy już komplet narzędzi: jądro kowariancji pozwala skonstruować
prawo procesu, formuły warunkowe opisują predykcję, a
oszacowania momentów dają regularność trajektorii.
Następny rozdział pokaże, że prosty wzór $K(s,t)=\min(s,t)$
prowadzi do procesu o bardzo bogatych, a nawet zaskakujących
własnościach geometrycznych.

% ==============================================================
% TOM II — ROZDZIAŁ 7: PROCES WIENERA
% ==============================================================
